\section{Planteamiento del problema}

El polimetilmetacrilato (PMMA), o acrílico, es un termoplástico amorfo con transición vítrea $T_g \approx$ \SI{105}{\celsius} \cite{gilormini2010,brandrup1999} y ventana de conformado por línea térmica entre \num{150} y \SI{160}{\celsius} \cite{plexiglas_forming,throne2008}. Su transparencia y facilidad de mecanizado lo hacen el material estructural más usado en montajes experimentales, carcasas de instrumentación y prototipos docentes.

El doblado por línea (\emph{line bending}) calienta una franja estrecha por encima de $T_g$ hasta el estado gomoso, la pliega y sostiene el ángulo mientras enfría \cite{osswald2012,throne2008}. La calidad depende de tres variables acopladas: el perfil de temperatura a lo ancho de la franja, que fija la zona plástica; el tiempo de permanencia, que gobierna la penetración térmica y la relajación de esfuerzos; y el enfriamiento bajo carga, que determina la recuperación elástica (\emph{springback}) y los esfuerzos residuales. Calentar de menos produce microfisuración (\emph{crazing}) y springback alto; calentar de más genera burbujas y deformación fuera de la línea \cite{yan2023,plexiglas_forming}. Enfriar bajo restricción, equivalente a un recocido corto, relaja los esfuerzos internos y reduce el springback \cite{osswald2012,yan2023}.

El estado del arte comercial presenta dos extremos. En la gama baja, los \emph{strip heaters} portátiles (anchos de \num{30}, \num{60} y \SI{125}{\centi\meter}, \SIrange{300}{1500}{\watt}, entre 80 y 300 USD) carecen de sensor y de realimentación: el operario ajusta por tiempo y experiencia y pliega a pulso contra plantilla \cite{shannon_line}. En la gama alta, las dobladoras de control numérico ofrecen calentamiento multilínea, prensado motorizado, ángulos de \SIrange{0}{180}{\degree} y refrigeración por agua, pero con longitudes útiles de \SIrange{1200}{3000}{\milli\meter} y precios de miles de dólares \cite{shannon_bending,mic_auto}. No se identificó oferta en el rango intermedio: equipos de mesa, de longitud útil inferior a \SI{300}{\milli\meter}, con control realimentado de temperatura y plegado automático. Ese es el nicho de un laboratorio universitario, donde las piezas son pequeñas y los lotes unitarios pero la repetibilidad importa igual que en producción. En el Laboratorio de Mecatrónica de EAFIT el doblado se hace hoy a mano sobre resistencias sin instrumentar: el resultado depende de la destreza del técnico, no es reproducible, consume tiempo calificado y genera rechazos.

\noindent\textbf{Pregunta de investigación:} \emph{¿Es posible diseñar, construir y validar una máquina de doblado por línea térmica de pequeño formato y bajo costo que, mediante control realimentado de la temperatura y plegado motorizado, produzca dobleces en PMMA de \SIrange{3}{6}{\milli\meter} de espesor y hasta \SI{25}{\centi\meter} de longitud con error angular no mayor a \SI{\pm 2}{\degree} y error de posición no mayor a \SI{\pm 1}{\milli\meter}?}
